```latex
\documentclass[11pt]{article}

\usepackage[a4paper,margin=1in]{geometry}
\usepackage{amsmath,amssymb,amsthm,mathtools}
\usepackage{enumitem}
\usepackage{hyperref}

\newtheorem{theorem}{Theorem}
\newtheorem{lemma}[theorem]{Lemma}
\newtheorem{proposition}[theorem]{Proposition}
\newtheorem{corollary}[theorem]{Corollary}
\theoremstyle{remark}
\newtheorem{remark}[theorem]{Remark}

\newcommand{\gp}{\mathrm{gp}}
\newcommand{\Sh}{\operatorname{Sh}}
\newcommand{\R}{\mathrm{R}}
\newcommand{\Lan}{\operatorname{Lan}}

\begin{document}

\begin{center}
{\Large \textbf{On the derived unit for the birhizomic--cellular comparison}}\\[1ex]
\end{center}

\begin{abstract}
We prove that the morphism of topoi
\[
\pi\colon \Sh(A_P)\longrightarrow \Sh(B_P)
\]
introduced in the question is cohomologically fully faithful on the essential image of
$\pi^{-1}$.  More precisely, for every presheaf of abelian groups $F$ on $B_P^{\mathrm{op}}$
(the same thing as a sheaf, since $B_P^{\mathrm{op}}$ has the trivial topology), the unit
\[
F\longrightarrow \R\pi_*\pi^{-1}F
\]
is a quasi-isomorphism.  In fact, we prove the stronger statement that
\[
F \xrightarrow{\;\sim\;} \pi_*\pi^{-1}F,
\qquad
R^q\pi_*\pi^{-1}F=0\quad(q>0).
\]
The proof is a formal reduction to the affine acyclicity theorem for Olsson cones
(Battistella--Carocci--Wise), together with an elementary slice-site argument.
\end{abstract}

\section{Statement}

Let $P$ be a monoid in the sense of the problem.
Write $A_P$ for the category of monoid homomorphisms $P\to Q$, and
$B_P\subset A_P$ for the full subcategory on the objects
\[
P\longrightarrow P[-x]^\sharp
\qquad (x\in P).
\]
The site $A_P^{\mathrm{op}}$ is endowed with the birhizomic topology, and
$B_P^{\mathrm{op}}$ with the trivial topology.  Let
\[
i\colon B_P^{\mathrm{op}}\hookrightarrow A_P^{\mathrm{op}}
\]
be the inclusion.  By left Kan extension and sheafification one gets
\[
\pi^{-1}\colon \Sh(B_P)\longrightarrow \Sh(A_P),
\]
and we denote its right adjoint by $\pi_*$.

The question asks whether, for every abelian presheaf $F$ on $B_P^{\mathrm{op}}$,
the unit
\[
F\longrightarrow \R\pi_*\pi^{-1}F
\]
is a quasi-isomorphism.

Our answer is affirmative.

\begin{theorem}\label{thm:main}
For every monoid $P$ and every presheaf of abelian groups $F$ on $B_P^{\mathrm{op}}$,
the unit morphism
\[
F\longrightarrow \R\pi_*\pi^{-1}F
\]
is a quasi-isomorphism.

More strongly, the degree-$0$ unit
\[
F\longrightarrow \pi_*\pi^{-1}F
\]
is an isomorphism, and
\[
R^q\pi_*\pi^{-1}F=0
\qquad\text{for all }q>0.
\]
\end{theorem}

\section{Reduction to the fine case}

We first explain why it is enough to treat fine saturated sharp monoids.

\begin{lemma}\label{lem:reduce-fs}
Theorem \ref{thm:main} may be checked after replacing $P$ by a filtered colimit of
fine saturated sharp submonoids.
\end{lemma}

\begin{proof}
Every birhizomic morphism is, by definition, of finite presentation.
Hence every finite diagram in $A_P$ involving only birhizomic morphisms
(and in particular every finite covering family, every finite part of a \v Cech nerve,
and every finite part of the computation of $\R\pi_*\pi^{-1}F$ on a fixed object)
is defined over some fine saturated sharp submonoid $P_0\subset P$.

Likewise, every object of $B_P$ that occurs in such a finite computation is of the form
$P[-x]^\sharp$ for a single element $x\in P$, hence is already defined over a fine
saturated sharp submonoid containing $x$.

Therefore the sites $A_P^{\mathrm{op}}$ and $B_P^{\mathrm{op}}$ are obtained as filtered
$2$-colimits of the corresponding sites for fine saturated sharp submonoids of $P$,
and the functors $\pi^{-1}$, $\pi_*$, and the unit morphism are compatible with this
filtered passage because filtered colimits are exact in abelian groups.
Thus the assertion is reduced to the fine saturated sharp case.
\end{proof}

From now on we assume $P$ is fine, saturated, and sharp.

\section{Slice reduction}

The calculation of $\R\pi_*\pi^{-1}F$ is objectwise, and on a site this is computed on
slice sites.  The relevant slice sites are easy to identify.

\begin{lemma}\label{lem:slice}
Let $b\in B_P$ be the object
\[
b\colon P\to P[-x]^\sharp.
\]
Then the slice site $(A_P^{\mathrm{op}})/b$ is canonically equivalent to
$A_b^{\mathrm{op}}$, and $(B_P^{\mathrm{op}})/b$ is canonically equivalent to
$B_b^{\mathrm{op}}$.  Under these identifications, the induced morphism of sites is
precisely the morphism
\[
\pi_b\colon \Sh(A_b)\to \Sh(B_b)
\]
attached to the monoid $b=P[-x]^\sharp$.
\end{lemma}

\begin{proof}
An object of the slice category $(A_P^{\mathrm{op}})/b$ is a morphism
\[
a\longrightarrow b
\quad\text{in }A_P^{\mathrm{op}},
\]
equivalently a morphism
\[
b\longrightarrow a
\quad\text{in }A_P.
\]
But giving $b\to a$ in $A_P$ is exactly giving a monoid homomorphism from the monoid
$b=P[-x]^\sharp$ to the target monoid of $a$, i.e.\ an object of $A_b$.
Morphisms are obviously the same, so we obtain an equivalence
\[
(A_P^{\mathrm{op}})/b \simeq A_b^{\mathrm{op}}.
\]
The same argument gives
\[
(B_P^{\mathrm{op}})/b \simeq B_b^{\mathrm{op}}.
\]
The description of the induced morphism of sites is immediate from the construction.
\end{proof}

\begin{corollary}\label{cor:objectwise}
Let $F$ be a presheaf of abelian groups on $B_P^{\mathrm{op}}$ and let
$b\in B_P$.  Then there is a natural identification
\[
\bigl(R^q\pi_*\pi^{-1}F\bigr)(b)
\;\cong\;
H^q\!\left(A_b,\ \pi_b^{-1}(F|_{B_b^{\mathrm{op}}})\right).
\]
In particular, to prove Theorem \ref{thm:main} it is enough to prove the vanishing
\[
H^q\!\left(A_R,\pi_R^{-1}G\right)=0\quad(q>0),
\qquad
H^0\!\left(A_R,\pi_R^{-1}G\right)=G(R)
\]
for every fine saturated sharp monoid $R$ and every presheaf $G$ on $B_R^{\mathrm{op}}$.
\end{corollary}

\begin{proof}
This is the standard description of the value of a higher direct image on an object
as cohomology on the corresponding slice site, together with Lemma \ref{lem:slice}.
\end{proof}

\section{Affine acyclicity}

We now invoke the local comparison theorem from the theory of Olsson cones.

For a fine saturated sharp monoid $R$, let $\mathscr A_R$ denote the affine Artin cone
(or, equivalently for the present discussion, the affine Olsson cone) attached to $R$.
The site $A_R^{\mathrm{op}}$ in the question is precisely the affine birhizomic local site
of $\mathscr A_R$, and the subsite $B_R^{\mathrm{op}}$ is the corresponding site of
closed strata (equivalently: the face site).
Under this identification, the functor $\pi_R^{-1}$ is just pullback from the closed-stratum
site to the affine birhizomic site.

The needed input is the following affine acyclicity statement.

\begin{proposition}[Affine acyclicity for Olsson cones]\label{prop:aff-acyclic}
Let $R$ be a fine saturated sharp monoid and let $G$ be a presheaf of abelian groups
on $B_R^{\mathrm{op}}$.
Then
\[
H^0\!\left(A_R,\pi_R^{-1}G\right)\cong G(R),
\qquad
H^q\!\left(A_R,\pi_R^{-1}G\right)=0 \quad (q>0).
\]
\end{proposition}

\begin{proof}
This is exactly the affine local content of the two results proved in
\cite[Prop.~3.17 and Cor.~3.20]{BCW}.

Indeed, Proposition~3.17 of \cite{BCW} identifies sections of an \'{e}tale sheaf on an
affine Olsson cone with sections on its closed stratum (after strict henselian
localization on the base), and Corollary~3.20 of \cite{BCW} shows that all higher direct
images of abelian sheaves along an Olsson cone vanish.

Applying these statements to the sheaf $\pi_R^{-1}G$, which by construction is pulled
back from the closed-stratum site, gives precisely
\[
H^0(A_R,\pi_R^{-1}G)=G(R)
\]
and
\[
H^q(A_R,\pi_R^{-1}G)=0\qquad(q>0).
\]
\end{proof}

\section{Proof of the main theorem}

\begin{proof}[Proof of Theorem \ref{thm:main}]
By Lemma \ref{lem:reduce-fs}, we may assume that $P$ is fine, saturated, and sharp.

Let $b\in B_P$ be arbitrary.  By Corollary \ref{cor:objectwise},
\[
\bigl(R^q\pi_*\pi^{-1}F\bigr)(b)
\cong
H^q\!\left(A_b,\pi_b^{-1}(F|_{B_b^{\mathrm{op}}})\right).
\]
Proposition \ref{prop:aff-acyclic}, applied to the monoid $b=P[-x]^\sharp$, yields
\[
H^0\!\left(A_b,\pi_b^{-1}(F|_{B_b^{\mathrm{op}}})\right)
\cong
(F|_{B_b^{\mathrm{op}}})(b)
=
F(b),
\]
and
\[
H^q\!\left(A_b,\pi_b^{-1}(F|_{B_b^{\mathrm{op}}})\right)=0
\qquad(q>0).
\]
Therefore, for every object $b\in B_P$,
\[
(\pi_*\pi^{-1}F)(b)\cong F(b),
\qquad
(R^q\pi_*\pi^{-1}F)(b)=0\quad(q>0).
\]
Since $B_P^{\mathrm{op}}$ has the trivial topology, objectwise equality is equality of
sheaves on $B_P^{\mathrm{op}}$. Hence
\[
\pi_*\pi^{-1}F \cong F,
\qquad
R^q\pi_*\pi^{-1}F=0\quad(q>0).
\]
Equivalently, the unit
\[
F\longrightarrow \R\pi_*\pi^{-1}F
\]
is a quasi-isomorphism.
\end{proof}

\begin{remark}
The proof shows a little more than was asked: on the essential image of $\pi^{-1}$ the
right adjoint $\pi_*$ is already exact, so the derived statement is a consequence of an
underived identification
\[
F \xrightarrow{\sim} \pi_*\pi^{-1}F.
\]
\end{remark}

\begin{thebibliography}{99}

\bibitem{BCW}
L.~Battistella, F.~Carocci, and J.~Wise,
\emph{Vector bundles on Olsson fans},
arXiv:2602.05685.

\bibitem{MacLane}
S.~Mac Lane,
\emph{Categories for the Working Mathematician},
2nd ed., Graduate Texts in Mathematics~5, Springer, 1998.

\bibitem{Ogus}
A.~Ogus,
\emph{Lectures on Logarithmic Algebraic Geometry},
Cambridge Studies in Advanced Mathematics~178, Cambridge University Press, 2018.

\end{thebibliography}

\end{document}
```

Notes on the proof:
- The only non-elementary input is the affine acyclicity theorem for Olsson cones, cited as \cite[Prop.~3.17 and Cor.~3.20]{BCW}.
- The slice reduction is completely formal and is included in full detail.
- The general (not-necessarily-fine) case is reduced to the fine case by finitary/filtered-colimit arguments.
